% Crossplane Reconciliation Loop — Chapter 5, Section 5.4.2
% Circular: XSDLC CR → Composition → Desired Children → Providers → Actual Resources → Status feedback
\begin{tikzpicture}[
    >={Stealth[length=2.5mm]},
    stage/.style={
        rectangle,
        rounded corners=4pt,
        draw=#1!60!black,
        fill=#1!12,
        minimum width=3cm,
        minimum height=0.9cm,
        font=\small,
        align=center
    },
    arrow/.style={
        ->,
        thick,
        #1!50!black
    }
]

% Circular layout (clockwise from top)
\node[stage=violet] (xsdlc)    at (0,3.5)    {\textbf{XSDLC CR}\\(desired state)};
\node[stage=teal]   (comp)     at (4.5,1.75)  {\textbf{Composition}\\(go-templating)};
\node[stage=blue]   (desired)  at (4.5,-1.75) {\textbf{Desired Children}\\(generated manifests)};
\node[stage=orange] (provider) at (0,-3.5)    {\textbf{Providers}\\(reconcile each child)};
\node[stage=green!60!black]  (actual)   at (-4.5,-1.75) {\textbf{Actual Resources}\\(cluster + GitHub)};
\node[stage=red!70!black]    (status)   at (-4.5,1.75)  {\textbf{Status Feedback}\\(ready/not-ready)};

% Clockwise arrows
\draw[arrow=violet]            (xsdlc.east)    -- node[above,font=\scriptsize,sloped] {spec input} (comp.north);
\draw[arrow=teal]              (comp.south)     -- node[right,font=\scriptsize] {render} (desired.north);
\draw[arrow=blue]              (desired.south)  -- node[below,font=\scriptsize,sloped] {create/update} (provider.east);
\draw[arrow=orange]            (provider.west)  -- node[below,font=\scriptsize,sloped] {provision} (actual.south);
\draw[arrow=green!60!black]    (actual.north)   -- node[left,font=\scriptsize] {observe} (status.south);
\draw[arrow=red!70!black,dashed] (status.north) -- node[above,font=\scriptsize,sloped] {conditions} (xsdlc.west);

% Center label
\node[font=\small\bfseries,text=gray!70,align=center] at (0,0) {Continuous\\Reconciliation};

% auto-ready annotation
\node[font=\tiny\itshape,text=red!50!black,align=center] at (-3.2,3.2) {function-auto-ready\\marks READY=True};

\end{tikzpicture}
